\section{Introduction}
La performance d'une application web détermine l'attente de l'utilisateur, le nombre de serveurs à louer, la
consommation d'énergie et la capacité à absorber un pic de trafic. Elle intéresse donc autant les développeurs que
les équipes qui exploitent le service, notamment dans les contextes où le matériel et la bande passante sont
limités.

Des affirmations circulent pourtant sans mesure précise : « Rust est plus rapide que PHP », « Node.js tient mieux la
charge que Python ». Les travaux publiés comparent surtout des langages sur des programmes de calcul isolés
\cite{pereira2017}, alors qu'une application web passe l'essentiel de son temps à lire du réseau, à interroger une
base de données et à produire du JSON. La méthodologie de la mesure compte autant que le résultat : variabilité
entre exécutions, moyennes trompeuses, latences extrêmes \cite{georges2007,dean2013}, et même la configuration du
test peut fausser les données \cite{mytkowicz2009}.

\textbf{Question de recherche.} Pour une même API de gestion de tâches adossée à SQLite, de combien les
performances de cinq technologies (Rust, Go, PHP, Node.js, Python) diffèrent-elles, et cet écart dépend-il de la
fonctionnalité sollicitée et du nombre d'utilisateurs simultanés ?

\textbf{Hypothèses.}
\textbf{H1} : sur les fonctionnalités légères (créer, lister, jointure simple), l'ordre est Rust, Go, puis PHP et
Node.js, et enfin Python, car le coût propre à la technologie (analyse HTTP, JSON, interpréteur) domine.
\textbf{H2} : sur les fonctionnalités où SQLite travaille longtemps (recherche, agrégation), les écarts se réduisent,
car le temps est passé dans le moteur de base de données, identique pour toutes les versions.
\textbf{H3} : la latence extrême (p95, p99) croît avec la charge, et plus vite pour les technologies les moins
performantes.
\textbf{H4} : les versions compilées (Rust, Go) consomment moins de mémoire et démarrent plus vite que les versions
interprétées.

\textbf{Contribution et plan.} Nous présentons cinq implémentations équivalentes d'une même API, un protocole de
mesure répété 30 fois par cas et les données brutes associées. La section~\ref{sec:methodo} décrit la méthode, la
section~\ref{sec:dev} la réalisation, la section~\ref{sec:res} les résultats et leur discussion, puis nous
concluons.
